All three defects reported by Prof.~I.~Grinis were root-caused and fixed with minimal, localised changes to \texttt{src/CategoryTheory.php} (and one dependent call site in \texttt{index.php}):

\begin{itemize}
  \item \texttt{Option::pure} now treats only \texttt{null} as absence of value, so \texttt{Some(false)} round-trips correctly through \texttt{bind(Option::pure)}, restoring the monad's right-identity law. The single call site that relied on the old behaviour (\texttt{fetchLiveJsonOption}, which used the \texttt{false} sentinel returned by \texttt{file\_get\_contents} on failure) was updated to test for that sentinel explicitly instead of overloading \texttt{Option::pure} with two meanings.
  \item \texttt{TripMetrics} now rounds to whole cents exactly once, at construction time, and stores that integer value (\texttt{costCents}). \texttt{append()} adds integers, which is exactly associative, instead of adding floats and re-rounding to two decimals after every step, which was not. The public float interface (\texttt{costEur}) is unchanged, so no other call site required modification.
  \item \texttt{Itinerary::advance()} now throws a \texttt{RuntimeException} carrying the free category's own \texttt{Left} error message when a graph edge does not compose with the current induced path, instead of silently keeping the previous (shorter, incorrect) path and returning a successful itinerary.
\end{itemize}

Each fix is accompanied by a new assertion in \texttt{tests/category\_laws\_test.php}, reproducing the exact counter-example supplied by Prof.~Grinis (including the $a=0.004,\,b=0.004,\,c=0.005$ monoid case). The three new assertions fail against the old code and pass against the fixed code, confirming both the regressions and the fixes. Separately, the project's additional invariant suite, \texttt{LietuvaTravel\_regression\_checks.php}, reports \textbf{10/10} checks passing against the corrected code. Table~\ref{tabl:summary} summarises the three defects end to end.

\begin{table}[!ht]\centering
\caption{Summary of the three defects, from report to fix to verification}
\label{tabl:summary}
\small
\begin{tabular}{|p{2.6cm}|p{3.6cm}|p{4.3cm}|p{3.3cm}|}\hline
\textbf{Law violated} & \textbf{Root cause} & \textbf{Fix} & \textbf{Verified by} \\\hline
Monad right identity ($m\mathbin{>\!\!>\!\!=}\mathrm{pure}\equiv m$) & \texttt{false} treated as absence in \texttt{Option::pure} & Only \texttt{null} denotes absence; sentinel checked explicitly at the one real caller & Unit assertion + live HTTP test (Sec.~\ref{sec:verification}) \\\hline
Monoid associativity ($(a\oplus b)\oplus c = a\oplus(b\oplus c)$) & Re-rounding to 2 decimals after every \texttt{append} & Accumulate in integer cents, round once at construction & Unit assertion with Prof.~Grinis' exact $a,b,c$ \\\hline
Error-propagation obligation of \texttt{Either::bind} & \texttt{Left} from \texttt{compose} read but not propagated in \texttt{advance()} & Throw on \texttt{Left}, carrying the categorical error message & Unit assertion asserting the throw occurs \\\hline
\end{tabular}
\end{table}

This progression -- a law-level description of what went wrong, a root-cause diagnosis in the code, a minimal fix, and an independent check that the fix actually holds -- was applied uniformly to all three defects, and is the same discipline the project's own law suite was supposed to enforce in the first place. The suite's earlier ``15/15'' badge was not dishonest, only incomplete: it verified every law it happened to test, but never generated the adversarial inputs (a falsy payload, a sub-cent amount, a deliberately broken graph) that expose these particular violations. The three new assertions close exactly that gap.
